\section{Searching for a closed form}
\label{a:sec:GA}

Several methods were tried in the hope of an elementary formula for $A(p)$. The
first was the continuum route of \cref{a:sec:nearcritical}: treat the map
as a differential equation in the large-generation limit and solve it. This gives
the shape of the decay but leaves the constant of integration undetermined, which
is to say it leaves $A(p)$ undetermined.

The second was symbolic regression by genetic algorithm. The idea is simple:
generate fifty candidate functions at random, measure how closely each reproduces
the target curve, retain the best five, and build a fresh population of fifty by
varying those; then measure again, and repeat the cycle of variation and selection
many times over \cite{stanley2002evolving,cohen2024physics}. The first
implementation searched a space of polynomials, on the hunch that the answer might
take that form and because such functions are simple to encode: a genome is a
coefficient list $[a_0,\dots,a_k]$ for $a_0+a_1p+\cdots+a_kp^k$ up to some maximum
order, drawn uniformly at random. It became apparent quickly that many
substantially different formulae produce curves indistinguishable from the target
by eye. Supposing that a true formula might have rational coefficients, a second
version searched $\sum_i(a_i/b_i)p^i$ with $[a_i],[b_i]\in\mathbb Z^{k+1}$, and a
third extended the space to ratios of two such polynomials. Candidates could be
constrained to the known endpoints $A(0)=\tfrac12$ and $A(\tfrac12)=0$, and the
best that emerged fitted the high-precision numerical data well:
\begin{equation}
\label{a:eq:A1hat}
\hat A_1(p) = \frac{\tfrac12(1-2p)(1+2p)}
{1+\tfrac12p-\tfrac52p^2-\tfrac65p^3+\tfrac65p^4} .
\end{equation}
This passes through $(0,\tfrac12)$ and $(\tfrac12,0)$ and looks convincing, but it
fails the diagnostic of \cref{a:rem:gradient}: differentiating at $p=\tfrac12$ gives
\begin{equation*}
\hat A_1'(\tfrac12) = -\frac{2}{0.55}\approx-3.636
\end{equation*}
rather than the required $-4$.

A later attempt used a symbolic tree search over progressively nested elementary
functions of all kinds. The conclusion was the same. Many functions of widely
varying form resemble the target closely, and no satisfactory simple formula
emerges even when candidates with fewer components are explicitly favoured. The
expressions returned were typically deep nestings of trigonometric and iterated
exponential terms carrying a dozen or more fitted decimal constants, and they were
not interpretable in any useful sense; a representative and comparatively concise
one is
\begin{equation}
\label{a:eq:A3hat}
\hat A_3(p) = \frac{-0.616647881739505\,\ee^{-0.552605335919693\,\ee^{p}}}
{\sin(\cos p)-|p|} + 0.921964323679771 .
\end{equation}
These are records of a curve-fitting exercise, not proposed identities.

\begin{figure}[ht]
\centering
\includegraphics[width=0.56\textwidth]{A3_hat_plot.pdf}
\caption{The symbolic-regression candidate $\hat A_3(p)$ of \cref{a:eq:A3hat}
(solid) against numerically computed values of $A(p)$ (red circles). The agreement
is excellent to the eye across the whole interval, which is exactly the problem:
the fit is not evidence of correctness. $\hat A_3$ misses the critical point,
returning $\hat A_3(\tfrac12) = 9.1\times10^{-4}$ rather than $0$, and arrives
there with gradient $-3.63$ instead of $-4$. Errors of this size are invisible here
but are fatal for $1/A$ near criticality.}
\label{a:fig:A3hat}
\end{figure}

The recurring defect is the same in every case. The candidates either fail to pass
through $(\tfrac12,0)$ or fail to arrive there with gradient $-4$, and by
\cref{a:rem:gradient} that is fatal near criticality.

A pragmatic patch is to splice the asymptotic onto a fitted curve,
\begin{equation}
\label{a:eq:A4hat}
\hat A_4(p) =
\begin{cases}
-0.04011483983619545\,\ee^{\ee^{2p}} + 0.6079994336528085, & p<0.499962,\\[2mm]
2-4p, & p\ge0.499962,
\end{cases}
\end{equation}
and this does work as a surrogate. But if one is going to splice, it is better to
splice onto something exact: the certified partial product
\cref{a:eq:producttailbound} on the first branch and the proved asymptotic
\cref{a:eq:Anearcritical} on the second. \Cref{a:sec:practical} does exactly that.

That an enormous number of very different formulae can behave almost identically over a bounded interval is not
merely an inconvenience of the method; it is a symptom. Curve fitting cannot
distinguish a function that has a closed form from one that does not, and no
quantity of fitting will ever settle the question.
